\chapter*{General Introduction}
\addcontentsline{toc}{chapter}{General Introduction} % to include the introduction to the table of content
\markboth{General Introduction}{} %To redefine the section page head

Online retail has become a central channel of commerce, yet in Tunisia
the digital marketplace landscape remains fragmented. Sellers who wish
to reach customers online are largely confined to two unsatisfying
options: vertically integrated single-vendor stores, which they cannot
join as independent merchants, and classified-ad sites, which list
offers but provide no order management, no quality control and no
buyer protection. The result is a market with weak tooling for content
moderation, product discovery and trust, precisely the qualities
that modern marketplaces are expected to guarantee.

It is in this context that the present project was carried out, as a
final-year engineering project within the company \textbf{NeuralBey}.
Its aim is to design and develop a \textbf{multi-vendor e-commerce
marketplace}: a platform on which independent sellers open and operate
their own stores within a shared environment governed by a central
operator. The platform lets sellers onboard with minimal friction,
publish products and fulfil orders, while administrators review seller
applications, moderate content and curate featured selections. Its
distinguishing capability is an \textbf{artificial-intelligence}
verification pipeline that automatically screens stores and products to
preserve buyer trust at scale, complemented by a conversational
database assistant for operators and an AI-assisted storefront redesign
for sellers. The system is engineered with a modern, fully
containerised technical foundation and an infrastructure-as-code
delivery pipeline.

To conduct the work in an organised and incremental manner, the project
adopted the agile \textbf{Scrum} method. This report is structured into
four chapters. The \textbf{first} presents the context and objectives
of the project: the host organisation, the study of existing solutions,
the proposed solution, the project-management methodology and the
modelling language. The \textbf{second} analyses and specifies the
requirements, actors, functional and non-functional requirements,
use-case and class diagrams, the Scrum project backlog and sprint
planning, the schedule, and the architecture and technical choices. The
\textbf{third} and \textbf{fourth} chapters describe the realisation of
the two successive releases, together with the tests and the
deployment. A general conclusion finally draws the assessment of the
work and its perspectives.
